%% back/colophon.tex -- hand-written, stable.
\chapter*{Colophon}
\hbtopicmark{COLOPHON}

\begingroup\small

\textbf{What this repository is.} A book whose body is a few hundred small
\LaTeX\ files, and whose build system is written so that adding a fourth,
fifth or fortieth source cannot damage the first three. The reading order is
generated from the metadata header of each file, so no index is ever
hand-maintained and no file is ever edited merely to make room for another.

\textbf{The two layers.} \texttt{topics/} holds the book: one idea per file,
fixed slot structure, curated synthesis. \texttt{sources/} holds the evidence:
one file per source per entry, append-only, never rewritten, with the
material found in only one source explicitly flagged. A flagged contribution
cannot be dropped from an entry --- the build fails if anyone tries. That is
the mechanism by which an early source survives a later one being better on
the same subject.

\textbf{Why the files are small.} Quality degrades with length, in human and
in machine writing alike. Repetition arrives, the voice drifts, and material
already written gets paraphrased back at itself. An entry that will not fit in
700 words is two entries, and splitting it is cheaper than repairing the
drift afterwards.

\textbf{The rules.} \texttt{docs/RULES.md} is the constitution and is
versioned. \texttt{AGENTS.md} is the short form for automated writers.
\texttt{tools/validate.py} enforces the machine-checkable subset;
\texttt{tools/churn\_guard.py} enforces the no-wholesale-rewrite rule against
git history. Nothing important is left to memory.

\textbf{State of the build.} This edition contains Parts I and II in full and
the scaffolding for Parts III to VII. The outstanding work is itemised in the
preceding appendix and in \texttt{registry/coverage.md}, both generated, so
neither can disagree with the book.

\textbf{Typesetting.} \texttt{style/handbook.cls} and \texttt{style/entry.sty}
hold every visual decision, and load each optional package behind a
conditional so that the book compiles on a bare \TeX\ Live installation. No
content file contains a layout command. To change how the book looks, change
those two files.

\par\endgroup
